\documentclass[10pt,a4paper]{article}
\usepackage[utf8]{inputenc}
\usepackage[T1]{fontenc}
\usepackage[catalan]{babel}
\usepackage{graphicx}
\usepackage{hyperref}
\usepackage[x11names]{xcolor}
\usepackage{fancyhdr}
\usepackage[margin=2cm]{geometry}
\usepackage{listings}
\graphicspath{{./imatges/}}% Lloc on hi ha els logos

\definecolor{backcolour}{rgb}{0.83, 0.83, 0.83}
\definecolor{textcolour}{rgb}{0.0, 0.0, 0.0}

\lstdefinestyle{mystyle}{
    backgroundcolor=\color{backcolour},  
    basicstyle=\ttfamily\footnotesize\color{textcolour}, 
    breakatwhitespace=false, 
    breaklines=true,  
    captionpos=b, 
    keepspaces=true,                 
    numbers=left, 
    numbersep=5pt, 
    showspaces=false, 
    showstringspaces=false, 
    showtabs=false, 
    tabsize=2
}

\lhead{Informe de la pràctica de Haskell}
\rhead{\thepage}
\cfoot{}
\begin{document}
    \begin{titlepage}
        \newcommand{\HRule}{\rule{\linewidth}{0.5mm}} % Defines a new command for the horizontal lines, change thickness here
        \begin{flushleft}
            \includegraphics[height=1.5cm]{EPS.png}\\\vfill
        \end{flushleft}
        \center % Center everything on the page
        %----------------------------------------------------------------------------------------
        %	HEADING SECTIONS
        %----------------------------------------------------------------------------------------
        \textsc{\huge \bfseries Paradigmes i llenguatges de programació }\\[0.25cm]
        \textsc{\Large \bfseries Curs 2024/25}\\[0.25cm]
        \textsc{\large GEINF }
        %----------------------------------------------------------------------------------------
        %	TITLE SECTION
        %----------------------------------------------------------------------------------------
        \HRule \\[0.4cm]
        { \huge \bfseries PAC4} \\[0.4cm] % Title of your document
{Juan José Gómez Villegas, \href{mailto:u1987338@campus.udg.edu}{u1987338@campus.udg.edu}} \\ {Guillem Pozo Sebastián, \href{mailto:u1972840@campus.udg.edu}{u1972840@campus.udg.edu}}
        \HRule \\\vfill
        \begin{minipage}{0.4\textwidth}
            \begin{flushleft}
		\footnotesize{Aquest document està subjecte a una llicència Creative Commons \href{https://creativecommons.org/licenses/by-nc-sa/4.0/?ref=chooser-v1}{CC BY-NC-SA 4.0}.} \\\vfill
		\includegraphics[height=1.5cm]{CC_license.png}
            \end{flushleft}
        \end{minipage}
        \hfill
        \begin{minipage}{0.4\textwidth}
            \begin{flushright} \large
                {\small (\today)}
            \end{flushright}
        \end{minipage}
    \end{titlepage}

\tableofcontents

\clearpage

\section{$\lambda$-càlcul clàssic}

\subsection{Tipus}

El tipus de dades \textbf{LT} ens permet representar $\lambda$-termes, i la definició és exactament la gramàtica que diu el que és un $\lambda$-terme, és a dir, un $\lambda$-terme és: una variable, una aplicació i una abstracció, res més és un $\lambda$-terme.

\lstset{language=Haskell, breaklines=true, basicstyle=\footnotesize}
\begin{lstlisting}[frame=mystyle]

data LT = Va String | Ap LT LT | Ab String LT

-- definim la forma de mostrar un lambda terme, els lambda termes es mostraran com (\x. (x y))
instance Show LT where
    show (Va a) = a
    show (Ap t1 t2) = "(" ++ show t1 ++ " " ++ show t2 ++ ")"
    show (Ab a t1) = "(\\" ++ a ++ ". " ++ show t1 ++ ")"

-- definim la idea d'equivalencia de dos lambda termes
instance Eq LT where
    (==) (Va a) (Va b) = a == b
    (==) (Ap t1 t2) (Ap t1' t2') = (&&) (t1 == t1') (t2 == t2')
    (==) (Ab _ t1) (Ab _ t2) = t1 == t2
    (==) _ _ = False
\end{lstlisting}

El \textbf{tipus LT} ha de ser instància de la classe \textbf{Show}, per tal que es pugui mostrar com nosaltres volguem, \textit{com diu a l'enunciat}, i ha de ser instància de la classe \textbf{Eq}, per tal de definir la idea d'equivalència de dos $\lambda$-termes, és a dir, dos $\lambda$-termes seran el mateix quan, siguin dues abstraccions i independentment de la variable, els seus $\lambda$-termes interns siguin iguals, quan siguin dues aplicacions i aquestes siguin iguals, quan siguin dos variables iguals. \\

També hem definit el tipus de dades \textbf{Substitució}, que \textbf{ens ha permès implementar la funció subst} tal com diu la teoria, és a dir, \textbf{com un operador que ens permet reemplaçar ocurrències de variables per termes, evitant captura de variables lliures}. Per tant, el tipus \textbf{Substitució} serà una variable $v$ que serà reemplaçada per un terme $M'$, sobre un terme $M$, definit a la funció \textbf{subst}, tot junt representarà: $M [v -> M']$.

\lstset{language=Haskell, breaklines=true, basicstyle=\footnotesize}
\begin{lstlisting}[frame=mystyle]
data Substitucio = Sub String LT
\end{lstlisting}

\clearpage

\subsection{Funcions auxiliars}

Les següents funcions ens permeten implementar les funcions principals com funcions d'ordre superior. \\

La funció \textbf{es\_redex} ens diu si un $\lambda$-terme té la forma d'un redex, o no.

\lstset{language=Haskell, breaklines=true, basicstyle=\footnotesize}
\begin{lstlisting}[frame=mystyle]
-- es_redex, funcio que donat un LT, retorna True si es un redex, False altrament
es_redex :: LT -> Bool
es_redex (Ap (Ab _ _) _) = True
es_redex _ = False
\end{lstlisting}

I per buscar redex en un $\lambda$-terme $T$, la funció \textbf{conte\_redex} ens diu si un terme qualsevol conté un redex a dins, \textbf{fent servir la funció es\_redex}, i funciona de la manera següent: Una variable no conté cap redex (retorna fals, cas base), una abstracció conté un redex, si i només si, el terme que conté a dins \textbf{és un redex}, i una aplicació conté un redex, si qualsevol dels dos termes de dins \textbf{són o contenen un redex}. Aquesta funció retorna cert si conté un redex, amb això ja podem simplificar el codi de les funcions principals.

\lstset{language=Haskell, breaklines=true, basicstyle=\footnotesize}
\begin{lstlisting}[frame=mystyle]
-- conte_redex, funcio que determina si un terme conte qualsevol redex (en qualsevol nivell, equivalent a fer esRedex sobre tot el terme)
conte_redex :: LT -> Bool
conte_redex (Va _) = False
conte_redex (Ab _ t) | es_redex t = True
                     | otherwise = conte_redex t
conte_redex (Ap t1 t2) | es_redex (Ap t1 t2) = True
                       | es_redex t1 = True
                       | es_redex t2 = True
                       | otherwise = conte_redex t1 || conte_redex t2
\end{lstlisting}

Les següents funcions ajuden a implementar la funció \textbf{subst}.

\lstset{language=Haskell, breaklines=true, basicstyle=\footnotesize}
\begin{lstlisting}[frame=mystyle]
-- is_var, funcio que diu si un LT es una variable
is_var :: LT -> Bool
is_var (Va _) = True
is_var _ = False

-- get_var, donat una variable d'un lambda terme retorna la variable com un string
get_var :: LT -> String
get_var (Va a) = a
get_var (Ab a _) = a
get_var (Ap t1 _) = get_var t1
\end{lstlisting}

I les següents funcions ajuden a implementar la funció auxiliar que retorna les variables lliures i lligades d'un terme qualsevol. La funció \textbf{eliminar\_duplicats}, donada una llista de variables, retorna la mateixa llista sense repetits, això soluciona un error trobar implementant la funció \textbf{freeAndboundVarsAux} en què les variables sortien repetides. I la funció \textbf{concat\_tuples} s'utilitza en la mateixa funció, en el cas que el terme entrat sigui una aplicació.

\lstset{language=Haskell, breaklines=true, basicstyle=\footnotesize}
\begin{lstlisting}[frame=mystyle]
-- eliminar_duplicats, funcio que elimina els elements duplicats d'una llista
eliminar_duplicats :: Eq a => [a] -> [a]
eliminar_duplicats = foldr (\x xs -> if x `elem` xs then xs else x:xs) []

-- concat_tuples, operador que concatena o intercalar dues llistes que son a dins d'una tupla
concat_tuples :: Eq a => ([a],[a]) -> ([a],[a]) -> ([a],[a])
concat_tuples t1 t2 = (eliminar_duplicats (fst t1 ++ fst t2),eliminar_duplicats (snd t1 ++ snd t2))

-- llargada, funcio que retorna la llargada d'una llista
llargada :: [a] -> Int
llargada = foldr (\_ y -> 1+y) 0
\end{lstlisting}

La següent funció a partir d'un $\lambda$-terme qualsevol, retorna una tupla amb dues llistes: primer una llista amb les variables lliures, segon una llista amb les variables lligades.

\lstset{language=Haskell, breaklines=true, basicstyle=\footnotesize}
\begin{lstlisting}[frame=mystyle]
-- freeAndboundVarsAux, funcio que construeix una tupla amb dues llistes que continguin les variables lliures (first) i lligades (second)
freeAndboundVarsAux :: [String] -> [String] -> LT -> ([String],[String])
freeAndboundVarsAux freeVars boundVars (Va a) | a `elem` boundVars = (freeVars,boundVars)
                                              | otherwise = (a:freeVars,boundVars)
freeAndboundVarsAux freeVars boundVars (Ab a t1) = freeAndboundVarsAux freeVars (a:boundVars) t1
freeAndboundVarsAux freeVars boundVars (Ap t1 t2) = freeAndboundVarsAux freeVars boundVars t1 `concat_tuples` freeAndboundVarsAux freeVars boundVars t2
\end{lstlisting}

\clearpage

\subsection{Funcions principals}

La funció \textbf{freeAndboundVars}, crida a la funció \textbf{freeAndboundVarsAux} com una funció d'ordre superior, inicialitzant dues llistes buides.

\lstset{language=Haskell, breaklines=true, basicstyle=\footnotesize}
\begin{lstlisting}[frame=mystyle]
-- freeAndboundVars, donat un LT retorna una tupla amb una llista de freeVars i una llista de boundVars
freeAndboundVars :: LT -> ([String],[String])
freeAndboundVars = freeAndboundVarsAux [] []
\end{lstlisting}

La subst \textbf{freeAndboundVars}, 

\lstset{language=Haskell, breaklines=true, basicstyle=\footnotesize}
\begin{lstlisting}[frame=mystyle]
-- subst, donat un LT i una Substitucio, retorna el mateix LT al que se li ha aplicat la Substitucio
subst :: LT -> Substitucio -> LT
subst (Va a) (Sub v t') | a == v = t'
                        | otherwise = Va a
subst (Ap t1 t2) (Sub v t') = Ap (subst t1 (Sub v t')) (subst t2 (Sub v t'))
subst (Ab a t1) (Sub v t') | a == v = if is_var t' then Ab (get_var t') (subst t1 (Sub v t')) else Ab a t1
                           | a /= v && a `notElem` fst (freeAndboundVars t') = Ab a (subst t1 (Sub v t'))
                           | a /= v && a `elem` fst (freeAndboundVars t') = subst (alfa_conv t1 t' a (a ++ "\'")) (Sub v t')
                           | otherwise = Ab a t1
    where alfa_conv t1 t' v' v'' | v'' `notElem` fst (freeAndboundVars t') && v'' `notElem` fst (freeAndboundVars t1) = Ab v'' (subst t1 (Sub v' (Va v'')))
                                 | otherwise = alfa_conv t1 t' v' (v'' ++ "\'")
\end{lstlisting}



\lstset{language=Haskell, breaklines=true, basicstyle=\footnotesize}
\begin{lstlisting}[frame=mystyle]
-- esta_normal, diu si LT ja esta en forma normal
esta_normal :: LT -> Bool
esta_normal = not . conte_redex
\end{lstlisting}



\lstset{language=Haskell, breaklines=true, basicstyle=\footnotesize}
\begin{lstlisting}[frame=mystyle]
-- beta_redueix, rep un LT que sigui un redex, i el resol
beta_redueix :: LT -> LT
beta_redueix (Ap (Ab a t1) t2) = subst t1 (Sub a t2)
beta_redueix (Ab _ t) = beta_redueix t
beta_redueix t = t
\end{lstlisting}



\lstset{language=Haskell, breaklines=true, basicstyle=\footnotesize}
\begin{lstlisting}[frame=mystyle]
-- redueix_un_n, rep un LT, i retorna el LT resultant d'aplicar la primera beta-reduccio segons l'ordre normal
redueix_un_n :: LT -> LT
redueix_un_n (Ap m n) | es_redex (Ap m n) = beta_redueix (Ap m n)
                      | conte_redex m = Ap (redueix_un_n m) n
                      | conte_redex n = Ap m (redueix_un_n n)
                      | otherwise = Ap m n
redueix_un_n (Ab x t) = Ab x (redueix_un_n t)
\end{lstlisting}



\lstset{language=Haskell, breaklines=true, basicstyle=\footnotesize}
\begin{lstlisting}[frame=mystyle]
-- redueix_un_a, rep un LT, i retorna el LT resultant d'aplicar la primera beta-reduccio segons l'ordre aplicatiu
redueix_un_a :: LT -> LT
redueix_un_a (Ap m n) | conte_redex m = Ap (redueix_un_a m) n
                      | conte_redex n = Ap m (redueix_un_a n)
                      | es_redex (Ap m n) = beta_redueix (Ap m n)
                      | otherwise = Ap m n
redueix_un_a (Ab x t) = Ab x (redueix_un_a t)
\end{lstlisting}



\lstset{language=Haskell, breaklines=true, basicstyle=\footnotesize}
\begin{lstlisting}[frame=mystyle]
-- l_normalitza_n, rep un LT, i retorna una llista de LT's que sigui una sequencia de beta-reduccions, segons l'ordre normal
l_normalitza_n :: LT -> [LT]
l_normalitza_n t | esta_normal t = t:[]
                 | otherwise = t:l_normalitza_n t'
    where t' = redueix_un_n t
\end{lstlisting}



\lstset{language=Haskell, breaklines=true, basicstyle=\footnotesize}
\begin{lstlisting}[frame=mystyle]
-- l_normalitza_a, rep un LT, i retorna una llista de LT's que sigui una sequencia de beta-reduccions, segons l'ordre aplicatiu
l_normalitza_a :: LT -> [LT]
l_normalitza_a t | esta_normal t = t:[]
                 | otherwise = t:l_normalitza_a t'
    where t' = redueix_un_a t
\end{lstlisting}



\lstset{language=Haskell, breaklines=true, basicstyle=\footnotesize}
\begin{lstlisting}[frame=mystyle]
-- normalitza_n, rep un LT, i retorna una tupla amb el nombre de passos, mes el LT en forma normal, seguint l'ordre normal
normalitza_n :: LT -> (Int, LT)
normalitza_n t = (llargada (lpassos t),last (lpassos t))
    where lpassos = l_normalitza_n
\end{lstlisting}



\lstset{language=Haskell, breaklines=true, basicstyle=\footnotesize}
\begin{lstlisting}[frame=mystyle]
-- normalitza_a, rep un LT, i retorna una tupla amb el nombre de passos, mes el LT en forma normal, seguint l'ordre aplicatiu
normalitza_a :: LT -> (Int, LT)
normalitza_a t = (llargada (lpassos t),last (lpassos t))
    where lpassos = l_normalitza_a
\end{lstlisting}

\clearpage

\section{Extra: Notació \textit{de Bruijn}}

\subsection{Tipus}

Per \textbf{la notació de Bruijn} hem definit un altre tipus de dades, que també és instància de les classes \textbf{Show} i \textbf{Eq}, igual que abans, l'únic que canvia en aquest tipus, és que ja no tenim variables en les $\lambda$-abstraccions, i les variables que abans eren cadenes de text, ara són enters, més concretament és la distància de cada variable en lambdes.

\begin{lstlisting}[language=Haskell]
data LTdB = VadB Int | ApdB LTdB LTdB | AbdB LTdB

-- que tambe sera instancia de la classe Show i Eq
instance Show LTdB where
    show (VadB a) = show a
    show (ApdB t1 t2) = "(" ++ show t1 ++ " " ++ show t2 ++ ")"
    show (AbdB t1) = "(\\" ++ ". " ++ show t1 ++ ")"

-- definim la mateixa idea d'equivalencia que teniem en els lambda termes pels lambda termes amb notacio de Bruijn
instance Eq LTdB where
    (==) (VadB a) (VadB b) = a == b
    (==) (ApdB t1 t2) (ApdB t1' t2') = (&&) (t1 == t1') (t2 == t2')
    (==) (AbdB t1) (AbdB t2) = t1 == t2
    (==) _ _ = False
\end{lstlisting}

També hem hagut de definir un Context.

\begin{lstlisting}[language=Haskell]
type Context = String
\end{lstlisting}

\clearpage

\subsection{Funcions auxiliars}

\clearpage

\subsection{Funcions principals}

\end{document}
